% === Begin Label Data ===
% ID: 300
% 难度星级: 
% 标签: 
% 备注: 
% 组卷引用次数: 0
% === End  Label Data ===

\begin{problem}{2026}{M}{2026年普通高等学校招生全国统一考试数学模拟测试(五)}{7}{函数}
已知函数 $ f(x) $ 的定义域为 $ \mathbb{Z} $，若 $ f(x+y) = f(x)f(1-y) + f(y)f(1-x) $，且 $ f(1) = -f(-1) = 1 $，则 $ \displaystyle\sum_{i=1}^{10}f(2i) = $
\begin{choices}
\choice{{$0$}}
\choice{{$1$}}
\choice{{$10$}}
\choice{{$20$}}
\end{choices}
\end{problem}

\begin{answer}
$0$
\end{answer}

\begin{solutions}
令 $ x = y = 0 $，得 $ f(0) = f(0)f(1) + f(0)f(1) = 2f(0)\cdot1 $，故 $ f(0) = 0 $.

令 $ x = y = 1 $，得 $ f(2) = f(1)f(0) + f(1)f(0) = 0 $.

令 $ y = 1 $，得 $ f(x+1) = f(x)f(0) + f(1)f(1-x) = f(1-x) $，即 $ f(x+1) = f(1-x) $，故函数关于直线 $ x = 1 $ 对称。

令 $ y = -1 $，得 $ f(x-1) = f(x)f(2) + f(-1)f(1-x) = 0 + (-1)f(1-x) = -f(1-x) $，结合上式 $ f(x+1) = f(1-x) $，得 $ f(x-1) = -f(x+1) $，即 $ f(t-2) = -f(t) $，故 $ f(t+4) = -f(t+2) = f(t) $，周期 $ T = 4 $.

由 $ f(0)=0 $，$ f(1)=1 $，$ f(2)=0 $，$ f(-1)=-1 $，得 $ f(3) = f(-1) = -1 $（由 $ f(x+1)=f(1-x) $，令 $ x=2 $ 得 $ f(3)=f(-1) $），验证 $ f(4)=f(0)=0 $.

故 $ f(2i) $：$ i=1 $ 时 $ f(2)=0 $；$ i=2 $ 时 $ f(4)=0 $；$ i=3 $ 时 $ f(6)=f(2)=0 $；… 每项均为 $ 0 $，所以和为 $ 0 $.
\end{solutions}